\subsection{Restitución funcional del momento orientado}
\label{iii:sec:restitucion-funcional-catalan}

La relación diferencial~\eqref{iii:eq:catalan-vacuum} admite una
inversa que conserva la colección completa de momentos. Su dominio es
la función constitutiva producida por los sectores del ciclo, con la
orientación de~\eqref{iii:eq:uv}. Las incidencias $90=3\cdot30$ y
$120=4\cdot30$ fijan esa función antes de evaluarla en los canales
$s_\pm=q_\pm^{30}$. La condición en el origen corresponde a la
desaparición del momento cuando se extingue su argumento contractivo.
Esta condición determina también las constantes de integración.

\begin{theorem}[Restitución integral y unicidad]
\label{iii:thm:restitucion-funcional-catalan}
Sea $f(s)=(1-s^3)/(1-s^4)$, para $0<s<1$, la respuesta
de~\eqref{iii:eq:trace-det}, y sea $\mathcal D=s\,d/ds$.
Existe una única función $F\in C^2((0,1))$ que satisface
\begin{equation}
 \mathcal D^2F(s)=\frac{s(1+s)}{1+s+s^2}f(s),
 \qquad \lim_{s\downarrow0}F(s)=0.
 \label{iii:eq:restitucion-funcional-problema}
\end{equation}
Está dada por
\begin{align}
 F(s)&=\int_0^s\log\frac{s}{t}\,
          \frac{1+t}{1+t+t^2}f(t)\,dt
       =\int_0^s\frac{\log(s/t)}{1+t^2}\,dt
       =\mathcal C_2(s),
 \label{iii:eq:restitucion-funcional-integral}\\
 f(s)&=\frac{1+s+s^2}{s(1+s)}\mathcal D^2F(s).
 \label{iii:eq:restitucion-funcional-inversa}
\end{align}
La extensión continua a $s=1$ recupera $F(1)=G_{\rm Cat}$.
\end{theorem}
\begin{proof}
La factorización $1+t+t^2+t^3=(1+t)(1+t^2)$ da
\[
 \frac{1+t}{1+t+t^2}f(t)=\frac1{1+t^2}.
\]
El integrando de~\eqref{iii:eq:restitucion-funcional-integral} es
integrable en cero, porque
\[
 0\leq F(s)\leq\int_0^s\log(s/t)\,dt=s.
\]
Por tanto $F(s)\to0$. La diferenciación en el intervalo abierto,
primero en un intervalo truncado y después por convergencia dominada,
produce
\[
 \mathcal DF(s)=\int_0^s\frac{dt}{1+t^2},
 \qquad \mathcal D^2F(s)=\frac{s}{1+s^2}
     =\frac{s(1+s)}{1+s+s^2}f(s).
\]
En particular, $F$ es de clase $C^2$ y satisface la ecuación.
Si $F_1,F_2$ son dos soluciones, la diferencia $H=F_1-F_2$
satisface $\mathcal D^2H=0$. Mediante $z=\log s$ se obtiene
$d^2H(e^z)/dz^2=0$, luego $H(s)=a\log s+b$.
La existencia del límite nulo cuando $s\downarrow0$ obliga primero
a $a=0$ y después a $b=0$. La condición sobre $F$ determina, pues,
la solución y también el límite $\mathcal DF(s)\to0$.

Para $s<1$, el desarrollo geométrico de $(1+t^2)^{-1}$ y la
integrabilidad absoluta de sus términos ponderados permiten integrar:
\[
 \int_0^s t^{2k}\log(s/t)\,dt
     =\frac{s^{2k+1}}{(2k+1)^2},\qquad
 F(s)=\sum_{k\geq0}\frac{(-1)^ks^{2k+1}}{(2k+1)^2}.
\]
Ésta es precisamente la colección~\eqref{iii:eq:c2}. La convergencia
absoluta y uniforme de esta última serie en $[0,1]$ autoriza el
paso al extremo $s=1$. También lo autoriza en la integral la cota
integrable $\log(1/t)$, extendiendo el integrando por cero para
$t>s$. Finalmente, despejar $f$ en la ecuación diferencial da
\eqref{iii:eq:restitucion-funcional-inversa}, con denominadores
positivos en $(0,1)$.
\end{proof}

El límite del momento y el límite de su derivada segunda se toman
en sus respectivas clases de convergencia. En particular,
\[
 \lim_{s\uparrow1}\mathcal D^2\mathcal C_2(s)=\frac12
\]
es el límite radial de Abel de la serie diferenciada. La evaluación
de $\mathcal C_2(1)$ utiliza directamente su serie absolutamente
convergente. Esta distinción conserva el mismo operador de profundidad
en el interior y en su lectura de frontera.

\subsection{Compatibilidad de la restitución con canales y velocidad}
\label{iii:sec:restitucion-canales-velocidad}

El teorema restituye una función a partir de otra función previamente
construida. La inversión cúbica~\eqref{iii:eq:cubic} actúa, en cambio,
sobre las evaluaciones $r_\pm$: recupera los argumentos $s_\pm$
dentro de esa colección fijada. La composición de ambas operaciones
conserva la orientación del cuarto de giro, las etiquetas $P_\pm$
y las incidencias del transporte. Así, $G_{\rm Cat}=\mathcal C_2(1)$
es la lectura extrema de un objeto funcional que participa también
en las respuestas interiores $r_\pm=f(s_\pm)$.

Para $\mathcal V_m=\mathcal V\otimes I_{9^m}$ y
$\tau_m=\operatorname{Tr}/(2\cdot9^m)$, la lectura de velocidad
mantiene exactamente la normalización de
\eqref{iii:eq:velocidad-traza-normalizada}:
\begin{equation}
 \exp[-2\tau_m(\log\mathcal V_m)]
   =\frac1{r_+r_-}
   =(\det\mathcal V)^{-1}=\widehat c.
 \label{iii:eq:restitucion-velocidad-normalizada}
\end{equation}
En efecto, cada autovalor $r_\pm$ se repite $9^m$ veces, de modo
que $2\tau_m(\log\mathcal V_m)=\log r_++\log r_-$.
La identidad utiliza el determinante de la realización de dos hojas
y la traza normalizada de su amplificación. Las secciones
dimensionales ya fijadas dan entonces
\[
 c_{\rm int}=c_{\rm vac}=\widehat c\,u_C,
 \qquad \ell_0=c_{\rm vac}t_0,
\]
con el mismo estado, la misma coordenatización y la misma duración elemental
de~\eqref{iii:eq:longitud-elemental-constitutiva}.
La restitución funcional conserva esta sección de velocidad al
componerla con la acción y la longitud; la multiplicidad espectral
queda absorbida por la normalización declarada.
